\begin{anexosenv}

\anexoitem{Referência Normativa -- Art. 6{\textordmasculine} da Lei Geral de Proteção de Dados Pessoais (Lei n{\textordmasculine} 13.709/2018)}{anexo:lgpd}

Transcrição do Art. 6{\textordmasculine} da Lei n{\textordmasculine} 13.709, de 14 de agosto de 2018 (LGPD) \cite{brasil2018lgpd}, utilizado como referência normativa para as restrições de privacidade e minimização de dados cadastrais no sistema \textbf{Karaokê Apache}:

\begin{citacao}
Art. 6{\textordmasculine} As atividades de tratamento de dados pessoais deverão observar a boa-fé e os seguintes princípios:

I -- finalidade: realização do tratamento para propósitos legítimos, específicos, explícitos e informados ao titular, sem possibilidade de tratamento posterior de forma incompatível com essas finalidades;

II -- adequação: compatibilidade do tratamento com as finalidades informadas ao titular, de acordo com o contexto do tratamento;

III -- necessidade: limitação do tratamento ao mínimo necessário para a realização de suas finalidades, com abrangência dos dados pertinentes, proporcionais e não excessivos em relação às finalidades do tratamento de dados;

VII -- segurança: utilização de medidas técnicas e administrativas aptas a proteger os dados pessoais de acessos não autorizados e de situações acidentais ou ilícitas de destruição, perda, alteração, comunicação ou difusão.
\end{citacao}

\end{anexosenv}
